\documentclass[12pt, letterpaper]{report}

\usepackage[top=2.5cm, bottom=2.5cm, left=3cm, right=2cm]{geometry}
\usepackage[spanish]{babel}
\usepackage{setspace}
\usepackage{amsmath,amssymb}
\usepackage{graphicx}
\usepackage{booktabs}
\usepackage{fancyhdr}
\usepackage{hyperref}
\usepackage{xcolor}

\definecolor{primary}{RGB}{255, 107, 53}
\onehalfspacing

\pagestyle{fancy}
\fancyhf{}
\fancyhead[R]{\small \slshape \nouppercase{\leftmark}}
\fancyfoot[C]{\thepage}

\begin{document}

% PORTADA OFICIAL DE TESIS
\begin{titlepage}
  \centering
  {\scshape\Large Universidad Nacional \par}
  {\scshape\large Departamento de Ciencias e Ingeniería \par}
  \vspace{3.5cm}

  {\Huge \bfseries \color{primary} DISEÑO E IMPLEMENTACIÓN DE UN ENTORNO LOCAL INTELIGENTE PARA LA EDICIÓN DE DOCUMENTOS CIENTÍFICOS \par}
  \vspace{1.5cm}
  {\Large Memoria de Título para optar al Grado Académico de\\ \textbf{Ingeniero Civil Informático} \par}
  \vspace{2.5cm}

  \begin{flushright}
    \large
    \textbf{Memorista:} Nombre del Estudiante \\
    \textbf{Profesor Guía:} Dr. Manuel E. Silva \\
    \textbf{Profesor Correferente:} Dr. Andrés C. Castro
  \end{flushright}

  \vfill
  {\large Valparaíso, Chile -- 2026 \par}
\end{titlepage}

% DEDICATORIA Y AGRADECIMIENTOS
\chapter*{Dedicatoria y Agradecimientos}
A mi familia, docentes y compañeros de ruta que hicieron posible este camino de investigación y desarrollo tecnológico. A la comunidad de código abierto por democratizar el conocimiento para las futuras generaciones.

% RESUMEN
\chapter*{Resumen}
En este trabajo de titulación se propone la arquitectura de un entorno de desarrollo integrado local orientado a la redacción de informes científicos en LaTeX. La solución elimina cuellos de botella de red y democratiza el trabajo en equipo en tiempo real mediante túneles seguros y asistencia por inteligencia artificial.

\tableofcontents
\listoffigures
\listoftables

% CAPÍTULOS
\chapter{Introducción}
\section{Contexto y Motivación}
La investigación académica contemporánea depende de la exactitud formal y tipográfica que ofrece el estándar LaTeX. Sin embargo, las alternativas disponibles en el mercado presentan barreras económicas o de conectividad.

\section{Objetivos}
\subsection{Objetivo General}
Diseñar, implementar y validar un entorno local ligero para la compilación y coedición de documentos académicos.

\chapter{Marco Teórico y Estado del Arte}
\section{Arquitecturas de Compilación TeX}
El compilador Tectonic transforma la distribución monolítica de TeX Live en un motor portable y autónomo capaz de gestionar dependencias dinámicamente.

\begin{equation}
  \mathcal{T}_{\text{build}} = f(\mathcal{C}_{\text{ast}}, \mathcal{F}_{\text{fonts}}, \mathcal{P}_{\text{packages}})
\end{equation}

\chapter{Metodología e Implementación}
\section{Módulos del Sistema}
El sistema se compone del núcleo de compilación en segundo plano, sincronización por WebSockets/HTTP y un asistente semántico basado en la API de Google Gemini.

\chapter{Resultados y Discusión}
Se comprobó la estabilidad del sistema con proyectos de más de 120 páginas, garantizando compilaciones en menos de 2 segundos.

\chapter{Conclusiones y Trabajo Futuro}
La implementación demostró que es posible alcanzar una experiencia fluida, robusta y gratuita sin depender de servidores en la nube.

\begin{thebibliography}{9}
\bibitem{knuth} D. Knuth, \textit{The TeXbook}, 1984.
\bibitem{tectonic} P. Williams, \textit{Tectonic Project Documentation}, 2023.
\end{thebibliography}

\end{document}
